\section{Fun Additive Results}
In this section we will see some fun results that concern or permeate
additive combinatorics.
Some of them may not be additive in nature but use tools common in
the additive world, so they may appear.

\section{The Combinatorial Nullstellensatz}
In this maybe wrongly located section we will explore nice applications
of a powerful theorem by Alon. The combinatorial Nullstellensatz
allows us to find non-zeros of certain polynomials with specific
degree condtions.
We will use it very much like in Section
	[\ref{section:extremal_set_theory}], devising
smart polynomials to get bounds on some nice problems.

\begin{theorem}[Alon's Combinatorial Nullstellensatz]
	\label{trm:combinatorial_nullstellensatz}
	Let $f \in \mathbb{F}[X_1, \dots, X_n]$ be a $n$ variable polynomial
	of degree $d = t_1 + \dots + t_n$,
	suppose furthermore the coefficient of $X_1^{t_1} \dots X_n^{t_n}$ is non zero.
	If there are sets $S_1, \dots S_n \subset \mathbb{F}$  with $|S_i| > t_i$,
	then there exists $X \in S_1 \times \dots \times S_n$ with $f(X) \neq 0$.
\end{theorem}

\begin{proof}
	We proceed by induction! For $n = 1$ the result is trivial, every
	polynomial $f$ of degree $d$ contains
	at most $d$ roots. Suppose it is valid for all numbers less or equal
	than $n - 1$, let's prove it for $n$.
	Let $f$ be given as in the assumption and consider the polynomial on
	one variable
	\[
		g_n(X_n) = \prod_{s \in S_n} (X_n - s) = X_N^{|S_n|} + \dots
	\]
	where the dots are lower order terms. Apply long division to $f$
	with $g_n$ to find $h$ and $r$ s.t
	\[
		f(X_1, \dots, X_n) =  h(X_1, \dots, X_n)g_n(X_n) + r(X_1, \dots, X_n)
	\]
	where crucially, $d(h) \leq d - |S_N|$ and the maximum degree of
	$X_n$ in $r$ is at most $|S_n| - 1$.
	So
	\[
		r(X_1, \dots, X_n) = \sum_{j = 0}^{|S_n| - 1} X_n^j r_{j}(X_1,
		\dots, X_{n-1}).
	\]
	Now, we know that the degree of $X_1^{t_1} \dots X_n^{t_n}$ of this
	thing is non zero. But, because
	$|S_n| > t_n$, it is clear that we can expand $h g_n$ as
	\[
		hg_n = h X_n^{|S_n|} + O(X_n^{d-1}).
	\]
	So $hg_n$ contains no monomial of type $X^\bm{t}$. In particular,
	it must be that $r_{t_n}$ has a non zero coefficient for $X_1^{t_1}
		\dots X_{n-1}^{t_{n-1}}$ and degree exactly $d$.
	So we can apply induction to it to find a non-zero root $x_1, \dots
		x_{n-1}$. Fixing this and considering
	$f(x_1, \dots, x_{n-1}, X_N)$ as a one variable polynomial ad
	applying the $n=1$ cases gives the result.
\end{proof}

Now let's see some applications of this thing. All of them are similar in the sense that we construct a smart polynomial
and apply this theorem, yielding a contradiction on the degree hypothesis.

\begin{theorem}

\end{theorem}

